\documentclass[11pt]{article}

\usepackage[T1]{fontenc}
\usepackage[utf8]{inputenc}
\usepackage{lmodern}
\usepackage[margin=1.08in]{geometry}
\usepackage{microtype}
\usepackage{amsmath,amssymb,amsthm,mathtools}
\usepackage{enumitem}
\usepackage{xcolor}
\usepackage{hyperref}
\usepackage[nameinlink,noabbrev]{cleveref}

\definecolor{linkblue}{RGB}{26,77,128}
\hypersetup{
  colorlinks=true,
  linkcolor=linkblue,
  citecolor=linkblue,
  urlcolor=linkblue,
  pdftitle={An adaptive packing bound for distinct consecutive sums},
  pdfauthor={AUTHOR NAME (REPLACE BEFORE RELEASE)},
  pdfsubject={A finite upper bound related to Erdos Problem 357},
  pdfkeywords={consecutive sums, additive combinatorics, Golomb rulers, Lean}
}

\setlist{nosep,leftmargin=2em}
\numberwithin{equation}{section}

\newtheorem{theorem}{Theorem}[section]
\newtheorem{proposition}[theorem]{Proposition}
\newtheorem{lemma}[theorem]{Lemma}
\newtheorem{corollary}[theorem]{Corollary}
\theoremstyle{remark}
\newtheorem{remark}[theorem]{Remark}

\newcommand{\N}{\mathbb{N}}
\newcommand{\ceil}[1]{\left\lceil #1\right\rceil}
\newcommand{\floor}[1]{\left\lfloor #1\right\rfloor}
\newcommand{\eps}{\varepsilon}

\title{\bfseries An adaptive packing bound\\
for distinct consecutive sums}
\author{AUTHOR NAME\\\small (replace before release)}
\date{Preprint draft, July 26, 2026}

\begin{document}
\maketitle

\begin{abstract}
Let \(f(n)\) be the largest \(k\) for which there are integers
\[
  1\leq a_1<\cdots<a_k\leq n
\]
such that all nonempty consecutive sums
\(\sum_{i=u}^{v}a_i\) are distinct.  Erd\H{o}s asked whether
\(f(n)=o(n)\).  We prove an adaptive finite inequality which, after
optimization, gives
\[
  f(n)\leq \frac n2+
  \left(\frac{9}{2^{7/3}}+o(1)\right)n^{2/3}.
\]
In particular, \(f(n)\leq (1/2+o(1))n\).  The proof attaches to each
starting value \(a_i\) a block of length
\(\ceil{n/(2a_i)}\).  Most of the resulting sums lie in an interval of
length \(n/2+O(n^{2/3})\); a telescoping estimate controls the starts
for which the block deviates too far from its initial value.  The
stated finite combinatorial inequality has also been encoded in
Lean~4, and the accompanying repository supplies an automated kernel
check.
\end{abstract}

\noindent\textbf{2020 Mathematics Subject Classification.}
11B75, 05D05.

\smallskip
\noindent\textbf{Keywords.}
Consecutive sums, additive combinatorics, Sidon sets, Golomb rulers,
formal verification.

\section{Introduction}

For a positive integer \(n\), let \(f(n)\) denote the maximum length of
a strictly increasing sequence
\[
  1\leq a_1<\cdots<a_k\leq n
\]
for which the sums
\begin{equation}
  A(u,v):=\sum_{i=u}^{v}a_i,
  \qquad 1\leq u\leq v\leq k,
  \label{eq:block-sums}
\end{equation}
are pairwise distinct.  Erd\H{o}s's 1977 paper asks a related
infinite-density question \cite{Erdos1977}.  The finite formulation
above is attributed to Erd\H{o}s and Harzheim and asks whether
\(f(n)=o(n)\); see Erd\H{o}s Problem~357
\cite{ErdosProblems357} and the problem collection of Erd\H{o}s and
Graham \cite{ErdosGraham1980}.

There is a useful additive reformulation.  Put \(s_0=0\) and
\(s_j=a_1+\cdots+a_j\).  The distinctness requirement says that all
positive differences \(s_v-s_u\), \(0\leq u<v\leq k\), are distinct.
Thus \(\{s_0,\ldots,s_k\}\) is a Golomb ruler, or a Sidon set in the
difference sense.  Moreover, its successive gaps \(a_1,\ldots,a_k\)
are strictly increasing, a configuration often called convex in
additive combinatorics.  The parameter \(n\) bounds the largest
successive gap rather than the diameter.

Known constructions give \cite{ErdosProblems357}
\[
  f(n)\geq (2+o(1))\sqrt n;
\]
computed initial values are recorded as OEIS A364132
\cite{OEISA364132}.  Thus the known upper and lower bounds remain far
apart.

Hegyv\'ari \cite{Hegyvari1986} studied the corresponding problem
without the monotonicity requirement and obtained a \(2n/3+o(n)\)
upper bound.  A theorem of Coppersmith and Phillips
\cite{CoppersmithPhillips1996}, proved under the weaker restriction
that no consecutive sum of at least two terms equals a later term,
implies
\begin{equation}
  f(n)\leq
  \left(\frac23-\frac1{512}\right)n+O(\log n).
  \label{eq:previous-bound}
\end{equation}
Indeed, pairwise distinct consecutive sums forbid such an equality.
More recently, Konieczny \cite{Konieczny2021} studied the number of
distinct consecutive sums in permutations, and Beker
\cite{Beker2024} studied the number of distinct values obtainable from
increasing sequences.  Those questions do not require every
consecutive sum to be different.

Our argument uses the strict increase of the \(a_i\)'s through an
adaptive choice of block length.  For a start \(i\), take
\begin{equation}
  r_i:=\ceil{\frac{n}{2a_i}}.
  \label{eq:adaptive-length}
\end{equation}
The main point is that \(r_i a_i\) always lies between \(n/2\) and
\(n\).  If a block of length \(r_i\) does not vary too much, its sum
therefore lies in a common interval of approximately \(n/2\) integer
values.  The lengths \(r_i\) are nonincreasing.  Within one fixed
length class, a telescoping argument bounds the \emph{total} variation
of all retained blocks, independently of the number of starts in the
class.

The resulting estimate still has leading term \(n/2\), so it does not
resolve the conjecture \(f(n)=o(n)\).  It does, however, improve the
previously available leading constant in \eqref{eq:previous-bound}.
More precisely, our main asymptotic consequence is
\[
  f(n)\leq \frac n2+\frac{9}{2^{7/3}}n^{2/3}+O(n^{1/3}).
\]
To the best of our knowledge, neither the adaptive finite inequality
below nor this leading constant has appeared previously.

\section{The finite adaptive bound}

We first state the parameterized form.  All divisions inside floor
symbols are ordinary real divisions.

\begin{theorem}[Finite adaptive inequality]
\label{thm:finite}
Let \(n,k,R,T\) be positive integers, and suppose that
\(1\leq a_1<\cdots<a_k\leq n\) has pairwise distinct nonempty
consecutive sums.  Then
\begin{align}
  k\leq {}&
  \floor{\frac{n+2T+2}{2}}
  +\floor{\frac{n+2R}{2R}}
  +\frac{R(R-1)}{2}
  +\floor{\frac{n(R-1)}{4T}}.
  \label{eq:finite-bound}
\end{align}
\end{theorem}

We prove the theorem by separating the starts into good starts and
three families of exceptions.

\subsection{Adaptive lengths and length classes}

For \(1\leq i\leq k\), define \(r_i\) by
\eqref{eq:adaptive-length}.  The elementary properties of the ceiling
function give
\begin{align}
  n&\leq 2a_i r_i, \label{eq:ceil-lower}\\
  2a_i(r_i-1)&<n \quad\text{if \(r_i>1\)}.
  \label{eq:ceil-upper}
\end{align}
We will also use
\begin{equation}
  r_i a_i\leq n.
  \label{eq:base-upper}
\end{equation}
For \(r_i=1\), this is just \(a_i\leq n\).  If \(r_i>1\), then
\(a_i\leq a_i(r_i-1)\), and hence
\[
  r_i a_i=a_i(r_i-1)+a_i
  \leq 2a_i(r_i-1)<n
\]
by \eqref{eq:ceil-upper}.  Since the \(a_i\)'s increase, the \(r_i\)'s
are nonincreasing.

Fix a cutoff \(R\).  Call a start \(i\) \emph{small-valued} if
\(r_i>R\), and denote the set of such starts by \(X\).
For \(i\in X\), the definition of the ceiling gives
\[
  2Ra_i<n.
\]
The values \(a_i\), \(i\in X\), are distinct positive integers and
each is at most
\[
  Q:=\floor{\frac{n+2R}{2R}}.
\]
Consequently
\begin{equation}
  |X|\leq \floor{\frac{n+2R}{2R}}.
  \label{eq:small-count}
\end{equation}
The deliberately loose extra \(2R\) in this estimate is retained to
match the integral formal statement.

For \(1\leq r\leq R\), let
\[
  C_r:=\{i:r_i=r\}.
\]
Because the \(r_i\)'s are nonincreasing, every nonempty \(C_r\) is an
interval of indices.  Write its first and last indices as
\(\ell_r\) and \(M_r\).  The boundary of \(C_r\) consists of those
starts \(i\in C_r\) for which \(i+r-1>M_r\).  There are at most
\(r-1\) such starts.  Summing over the classes gives
\begin{equation}
  \#\{\text{boundary starts}\}
  \leq \sum_{r=1}^{R}(r-1)
  =\frac{R(R-1)}2.
  \label{eq:boundary-count}
\end{equation}

The remaining starts in \(C_r\) form its interior:
\[
  I_r:=\{i\in C_r:i+r-1\leq M_r\}.
\]
If \(i\in I_r\), then every index \(i,i+1,\ldots,i+r-1\)
belongs to \(C_r\).  Indeed, for \(i\leq j\leq M_r\), monotonicity
gives \(r_{M_r}\leq r_j\leq r_i\), and both endpoints equal \(r\).
In particular, the length-\(r\) block beginning at an interior start
is contained in the original sequence.

\subsection{The class error estimate}

For \(i\in I_r\), define the block error
\begin{equation}
  E_i^{(r)}
  :=\sum_{t=0}^{r-1}(a_{i+t}-a_i).
  \label{eq:error}
\end{equation}
Thus the corresponding block sum is
\begin{equation}
  \sum_{t=0}^{r-1}a_{i+t}
  =r a_i+E_i^{(r)}.
  \label{eq:block-decomposition}
\end{equation}

\begin{lemma}[Telescoping within one class]
\label{lem:class-error}
For every \(r\geq1\),
\begin{equation}
  4\sum_{i\in I_r}E_i^{(r)}\leq n.
  \label{eq:class-error-bound}
\end{equation}
\end{lemma}

\begin{proof}
There is nothing to prove if \(I_r\) is empty.  Otherwise set
\[
  \ell=\ell_r,\qquad M=M_r,\qquad
  L=a_\ell,\qquad U=a_M.
\]
Then \(I_r=\{\ell,\ell+1,\ldots,M-r+1\}\).
For \(0\leq t\leq r-1\), cancellation of the common middle terms
gives
\begin{align*}
  \sum_{i=\ell}^{M-r+1}(a_{i+t}-a_i)
  &=
  \sum_{j=M-r+2}^{M-r+1+t}a_j
  -\sum_{j=\ell}^{\ell+t-1}a_j\\
  &\leq t(U-L).
\end{align*}
Here an empty sum is interpreted as zero.  Summing in \(t\), and
interchanging the two finite sums, yields
\begin{equation}
  \sum_{i\in I_r}E_i^{(r)}
  \leq \frac{r(r-1)}2\,(U-L).
  \label{eq:error-before-width}
\end{equation}

It remains to control the width \(U-L\) of the class.  If \(r=1\),
the right-hand side of \eqref{eq:error-before-width} is zero.  Assume
\(r>1\).  Applying \eqref{eq:ceil-lower} at \(\ell\) and
\eqref{eq:ceil-upper} at \(M\) gives
\[
  n\leq2Lr,\qquad 2U(r-1)<n.
\]
After multiplying the second inequality by \(r\), and the first by
\(r-1\), we obtain
\[
  2r(r-1)(U-L)<nr-n(r-1)=n.
\]
Combining this with \eqref{eq:error-before-width} proves
\eqref{eq:class-error-bound}.
\end{proof}

Call an interior start \(i\in I_r\) \emph{bad} if
\(E_i^{(r)}>T\).  Let \(D_r\) denote the bad starts in class \(r\).
Since all errors are nonnegative, \cref{lem:class-error} gives
\[
  4T|D_r|
  \leq4\sum_{i\in D_r}E_i^{(r)}
  \leq4\sum_{i\in I_r}E_i^{(r)}
  \leq n.
\]
For \(r=1\) the error is identically zero, so \(D_1\) is empty.
The length classes are disjoint.  Summing the preceding inequality
over \(r=2,\ldots,R\) gives
\[
  4T\,\#\{\text{bad starts}\}\leq n(R-1),
\]
and therefore
\begin{equation}
  \#\{\text{bad starts}\}
  \leq
  \floor{\frac{n(R-1)}{4T}}.
  \label{eq:bad-count}
\end{equation}

\subsection{Packing the good block sums}

Every start not already counted in \eqref{eq:small-count},
\eqref{eq:boundary-count}, or \eqref{eq:bad-count} is called
\emph{good}.  If \(i\) is good, its block length is \(r_i\), its block
is contained in the sequence, and its error is at most \(T\).
Equations \eqref{eq:ceil-lower} and \eqref{eq:base-upper} imply
\[
  \ceil{\frac n2}\leq r_i a_i\leq n.
\]
Together with \eqref{eq:block-decomposition}, this gives
\begin{equation}
  \ceil{\frac n2}
  \leq \sum_{t=0}^{r_i-1}a_{i+t}
  \leq n+T.
  \label{eq:good-interval}
\end{equation}

The blocks attached to distinct good starts are distinct blocks, and
all consecutive block sums were assumed pairwise distinct.
Consequently, the sums in \eqref{eq:good-interval} are distinct
integers.  The interval there contains
\[
  n+T-\ceil{\frac n2}+1
  =\floor{\frac{n+2T+2}{2}}
\]
integers.  Thus
\begin{equation}
  \#\{\text{good starts}\}
  \leq\floor{\frac{n+2T+2}{2}}.
  \label{eq:good-count}
\end{equation}

Every start is good, small-valued, a class-boundary start, or bad.
Adding \eqref{eq:small-count}, \eqref{eq:boundary-count},
\eqref{eq:bad-count}, and \eqref{eq:good-count} proves
\cref{thm:finite}.

\section{Consequences and optimization}

Applying \cref{thm:finite} to a sequence of maximum length immediately
gives the same inequality with \(k\) replaced by \(f(n)\).

\begin{corollary}[A simple cube-scale bound]
\label{cor:cube}
If \(m\geq1\) and \(n\leq m^3\), then
\begin{equation}
  f(n)\leq \frac n2+\frac94m^2+2.
  \label{eq:simple-cube}
\end{equation}
In particular, for \(m=\ceil{n^{1/3}}\),
\[
  f(n)\leq\frac n2+\frac94n^{2/3}+O(n^{1/3}).
\]
\end{corollary}

\begin{proof}
Take \(R=m\) and \(T=m^2\) in \cref{thm:finite}, discard the
floors, and use \(n\leq m^3\).  The four contributions beyond \(n/2\)
are at most
\[
  m^2+1,\qquad
  \frac{m^2}{2}+1,\qquad
  \frac{m^2}{2},\qquad
  \frac{m^2}{4},
\]
respectively.
\end{proof}

The two free parameters permit a better constant.

\begin{corollary}[Optimized asymptotic form]
\label{cor:optimized}
As \(n\to\infty\),
\begin{equation}
  f(n)\leq\frac n2+
  \frac{9}{2^{7/3}}\,n^{2/3}
  +O(n^{1/3}).
  \label{eq:optimized}
\end{equation}
In particular,
\[
  f(n)\leq\left(\frac12+o(1)\right)n.
\]
\end{corollary}

\begin{proof}
For fixed \(c,t>0\), take
\[
  R=\ceil{c\,n^{1/3}},
  \qquad
  T=\ceil{t\,n^{2/3}}.
\]
The finite inequality gives
\begin{align*}
  f(n)
  \leq \frac n2+
  \left(
    t+\frac1{2c}+\frac{c^2}{2}+\frac{c}{4t}
  \right)n^{2/3}
  +O(n^{1/3}).
\end{align*}
For fixed \(c\), the expression is minimized in \(t\) at
\(t=\sqrt{c}/2\).  It then becomes
\[
  \sqrt{c}+\frac1{2c}+\frac{c^2}{2}.
\]
Writing \(y=c^{3/2}\), the critical-point equation is
\[
  2y^2+y-1=0,
\]
whose positive solution is \(y=1/2\).  Thus
\[
  c=2^{-2/3},\qquad t=2^{-4/3},
\]
and the resulting coefficient is
\[
  2^{-1/3}+2^{-1/3}+2^{-7/3}
  =\frac{9}{2^{7/3}}.
\]
The objective tends to infinity as \(c\) tends to either \(0\) or
\(\infty\), so this critical point gives the global minimum.
\end{proof}

\begin{remark}
Numerically, \(9/2^{7/3}\approx1.7858\).  The secondary-term
optimization does not change the central obstruction: the good block
sums occupy an interval with asymptotic length \(n/2\).  Within this
packing framework, a proof of \(f(n)=o(n)\) would therefore need an
additional idea that substantially reduces, or exploits the
distribution of, those good sums.
\end{remark}

\section{Formal development and verification protocol}

The finite inequality \eqref{eq:finite-bound} was formalized in
Lean~4 using natural-number arithmetic throughout.  The formal
statement uses a function \(a:\operatorname{Fin}(k)\to\mathbb Z\),
the range condition \(a_i\in[1,n]\), strict monotonicity, and
injectivity of the sum map on order-connected finite subsets.  These
are the same data used in the Formal Conjectures formulation of
Erd\H{o}s Problem~357 \cite{FormalConjectures}.

The formal development constructs four finite sets corresponding to
the good, small-valued, boundary, and bad starts.  Its key
denominator-free estimates are
\begin{align*}
  2|G|&\leq n+2T+2,\\
  2R|X|&\leq n+2R,\\
  2|B|&\leq R(R-1),\\
  4T|D|&\leq n(R-1).
\end{align*}
It then derives \eqref{eq:finite-bound} by integer division.  As a
convenient fully integral specialization, the formal file also proves
that, whenever \(n\leq m^3\),
\begin{equation}
  2f(n)\leq n+8m^2+4.
  \label{eq:lean-cube}
\end{equation}
This last displayed estimate is deliberately coarser than
\cref{cor:cube}; its role is to provide a short exact certificate of
the \(n/2+O(n^{2/3})\) scale.

The accompanying source file is
\texttt{Erdos357AdaptiveBound.lean}.  It contains 1841 lines, and a
static source audit finds no proof placeholder, user-declared axiom, or
unsafe declaration.  The repository pins Lean~4.27.0 and
mathlib~v4.27.0.  Its continuous-integration workflow builds the
project and invokes an additional checker; the repository's
\texttt{VERIFICATION.md} records the source digest and the distinction
between static, typesetting, and kernel checks.  A public repository
URL and a green continuous-integration run should be cited in the
released version.

\section{Concluding remarks}

The argument is driven by two elementary features that work together:
the adaptive base term \(r_i a_i\) is automatically in
\([\ceil{n/2},n]\), while the total block error in a fixed length class
telescopes to a bound depending only on the width of that class.
The relation
\[
  2r(r-1)(U-L)<n
\]
is precisely what cancels the quadratic triangular factor produced by
the block errors.

Several natural questions remain.  Can one show that the good sums
avoid a positive proportion of
\([\ceil{n/2},n+T]\)?  Can information from different adaptive length
classes be combined, rather than treating their sums only as a single
injective family?  Finally, can the convex-Sidon interpretation supply
additional restrictions on how these differences are distributed?
Any improvement of the leading \(1/2\) appears to require progress on
one of these points.

\section*{Disclosure}

This manuscript and its repository documentation were drafted and
edited with assistance from OpenAI language models, starting from the
supplied Lean development.  The named human author or authors are
responsible for independently checking the mathematical exposition and
for every submitted claim.  The Lean kernel and the additional checker,
not a language model, provide the machine-verification steps.  This
disclosure should be revised, if necessary, to comply with the selected
journal's current policy.

\begin{thebibliography}{99}

\bibitem{Beker2024}
A.~Beker,
\newblock On a problem of Erd\H{o}s and Graham about consecutive sums
in strictly increasing sequences,
\newblock \emph{Bull. Lond. Math. Soc.} \textbf{56} (2024), no.~8,
2749--2759.
\newblock \url{https://doi.org/10.1112/blms.13098}.

\bibitem{CoppersmithPhillips1996}
D.~Coppersmith and S.~Phillips,
\newblock On a question of Erd\H{o}s on subsequence sums,
\newblock \emph{SIAM J. Discrete Math.} \textbf{9} (1996), no.~2,
173--177.
\newblock \url{https://doi.org/10.1137/S0895480193244139}.

\bibitem{Erdos1977}
P.~Erd\H{o}s,
\newblock Problems and results on combinatorial number theory III,
\newblock in \emph{Number Theory Day}, Lecture Notes in Mathematics
626, Springer, Berlin, 1977, pp.~43--72.
\newblock \url{https://www.renyi.hu/~p_erdos/1977-27.pdf}.

\bibitem{ErdosProblems357}
T.~Bloom and T.~Tao (maintainers),
\newblock Erd\H{o}s Problem 357,
\newblock \emph{Erd\H{o}s Problems database}, accessed July 26, 2026.
\newblock \url{https://www.erdosproblems.com/357}.

\bibitem{ErdosGraham1980}
P.~Erd\H{o}s and R.~L. Graham,
\newblock \emph{Old and New Problems and Results in Combinatorial
Number Theory},
\newblock Monographies de L'Enseignement Math\'ematique 28,
Universit\'e de Gen\`eve, 1980.

\bibitem{FormalConjectures}
M.~Firsching, P.~Lezeau, S.~Mercuri, et al.,
\newblock Formal Conjectures: an open and evolving benchmark for
verified discovery in mathematics,
\newblock arXiv:2605.13171, 2026.
\newblock See the
\href{https://github.com/google-deepmind/formal-conjectures/blob/c252a41054125b5fd9c8356e2137cd9b55337657/FormalConjectures/ErdosProblems/357.lean}
{Problem~357 source pinned at commit \texttt{c252a410}}.

\bibitem{Hegyvari1986}
N.~Hegyv\'ari,
\newblock On consecutive sums in sequences,
\newblock \emph{Acta Math. Hungar.} \textbf{48} (1986), no.~1--2,
193--200.
\newblock \url{https://doi.org/10.1007/BF01949064}.

\bibitem{Konieczny2021}
J.~Konieczny,
\newblock On consecutive sums in permutations,
\newblock \emph{J. Comb.} \textbf{12} (2021), no.~3, 413--477.
\newblock \url{https://doi.org/10.4310/JOC.2021.v12.n3.a3}.

\bibitem{OEISA364132}
OEIS Foundation Inc.,
\newblock Entry A364132,
\newblock \emph{The On-Line Encyclopedia of Integer Sequences},
accessed July 26, 2026.
\newblock \url{https://oeis.org/A364132}.

\end{thebibliography}

\end{document}
